% FORMAT AND PACKAGES
% {
\documentclass[a4paper]{article}
\usepackage{a4wide,amssymb,epsfig,latexsym,multicol,array,hhline,fancyhdr}
\usepackage{tcolorbox}
\usepackage{minted}
\usepackage{vntex}
\usepackage{amsmath}
\usepackage{lastpage}
\usepackage[lined,boxed,commentsnumbered]{algorithm2e}
\usepackage{enumerate}
\usepackage{xcolor}
\usepackage{graphicx}							% Standard graphics package
\usepackage{array}
\usepackage{tabularx, caption}
\usepackage{multirow}
\usepackage{multicol}
\usepackage{rotating}
\usepackage{graphics}
\usepackage{geometry}
\usepackage{setspace}
\usepackage{epsfig}
\usepackage{tikz}
\usepackage{xfrac}
\usepackage{bm}
\usepackage{biblatex}
\usepackage[colorlinks]{hyperref}
\newcommand{\cach}{\hspace*{1.5em}\ignorespaces}
% \usepackage[acronym,toc]{glossaries}
% \usepackage[symbols,nogroupskip,nonumberlist]{glossaries-extra}
\usepackage[
 sort=none,% no sorting or indexing required
 abbreviations,% create list of abbreviations
 symbols,% create list of symbols
 stylemods,style=list, % set the default glossary style
 nogroupskip, nonumberlist, nomain
]{glossaries-extra}


% FORMATTING
% {
\DeclareMathOperator{\arccot}{arccot}
\captionsetup[table]{name=Bảng}
\captionsetup[figure]{name=Hình}
\newenvironment{Description}{\list{}{%
    \let\makelabel\descriptionlabel    % this comes from the original description environment
    \setlength{\rightmargin}{\leftmargin}% this comes from the original quote environment
    \setlength{\labelwidth}{0pt}%          this is new
    }}{\endlist}

\addbibresource{citations.bib}
    
\hypersetup{urlcolor=blue,linkcolor=black,citecolor=black,colorlinks=true} 
\usetikzlibrary{arrows,snakes,backgrounds}
\definecolor{mathblue}{RGB}{0,114,188}
% \makeatletter  \def\m@th{\mathsurround\z@\color{mathblue}} \makeatother
% \everymath{\color{mathblue}}
% \setmathfont[Color=000000]{Arial}
%\usepackage{pstcol} 								% PSTricks with the standard color package
\newtheorem{theorem}{{\bf Theorem}}
\newtheorem{property}{{\bf Property}}
\newtheorem{proposition}{{\bf Proposition}}
\newtheorem{corollary}[proposition]{{\bf Corollary}}
\newtheorem{lemma}[proposition]{{\bf Lemma}}

\AtBeginDocument{\renewcommand{\listfigurename}{Danh sách hình ảnh}}
\AtBeginDocument{\renewcommand{\listtablename}{Danh sách bảng biểu}}
\AtBeginDocument{\renewcommand*\contentsname{Mục lục}}
\AtBeginDocument{\renewcommand*\refname{Tài liệu tham khảo}}
%\usepackage{fancyhdr}

\setlength{\headheight}{40pt}
\pagestyle{fancy}
\fancyhead{} % clear all header fields
\fancyhead[L]{
 \begin{tabular}{rl}
    \begin{picture}(25,15)(0,0)
    \put(0,-8){\includegraphics[width=8mm, height=8mm]{hcmut.png}}
    %\put(0,-8){\epsfig{width=10mm,figure=hcmut.eps}}
   \end{picture}&
	%\includegraphics[width=8mm, height=8mm]{hcmut.png} & %
	\begin{tabular}{l}
		\textbf{\bf \ttfamily Trường Đại học Bách Khoa TP.Hồ Chí Minh}\\
		\textbf{\bf \ttfamily Khoa Khoa học và Kỹ thuật máy tính}
	\end{tabular} 	
 \end{tabular}
}
\fancyhead[R]{
	\begin{tabular}{l}
		\tiny \bf \\
		\tiny \bf 
	\end{tabular}  }
\fancyfoot{} % clear all footer fields
\fancyfoot[L]{\scriptsize \ttfamily Báo cáo Bài tập lớn Đồ án tổng hợp CNPM - HK251 - Năm học 2025 - 2026}
\fancyfoot[R]{\scriptsize \ttfamily Trang {\thepage}/\pageref{LastPage}}
\renewcommand{\headrulewidth}{0.3pt}
\renewcommand{\footrulewidth}{0.3pt}

\setcounter{secnumdepth}{4}
\setcounter{tocdepth}{4}

\makeatletter
\newcounter {subsubsubsection}[subsubsection]
\renewcommand\thesubsubsubsection{\thesubsubsection .\@alph\c@subsubsubsection}
\newcommand\subsubsubsection{\@startsection{subsubsubsection}{4}{\z@}%
                                     {-3.25ex\@plus -1ex \@minus -.2ex}%
                                     {1.5ex \@plus .2ex}%
                                     {\normalfont\normalsize\bfseries}}
\newcommand*\l@subsubsubsection{\@dottedtocline{3}{10.0em}{4.1em}}
\newcommand*{\subsubsubsectionmark}[1]{}
% \def\m@th{\mathsurround\z@\color{mathblue}}
\makeatother
% }
% }

% ACRONYMS & SYMBOLS
% {
% \makeglossaries
\setabbreviationstyle{long-short}
\newabbreviation{csp}{CSP}{Cutting Stock Problem}
\newabbreviation{ffd}{FFD}{First Fit Decreasing}
\newabbreviation{ga}{GA}{Genetic Algorithm}
\newabbreviation{lp}{LP}{Linear Programming}
% \glsnoexpandfields
\glsxtrnewsymbol[description = {Tập hợp số tự nhiên}]{natural}{\ensuremath{\mathbb{N}}}

% }
%
% DOCUMENT
\begin{document}

% TITLE PAGE
\begin{titlepage}
\begin{center}
ĐẠI HỌC QUỐC GIA THÀNH PHỐ HỒ CHÍ MINH\\
TRƯỜNG ĐẠI HỌC BÁCH KHOA\\
KHOA KHOA HỌC VÀ KỸ THUẬT MÁY TÍNH\\
\end{center}

\vspace{1cm}

\begin{figure}[h!]
\begin{center}
\includegraphics[width=3cm]{hcmut.png}
\end{center}
\end{figure}

\vspace{1cm}


\begin{center}
\begin{tabular}{c}
\multicolumn{1}{c}{\textbf{{\Large Đồ án tổng hợp - CNPM (CO3103)}}}\\
~~\\
\hline
\\
\multicolumn{1}{l}{\textbf{{\Large Bài tập lớn}}}\\
\\
\textbf{\textit{{\Huge "Ứng dụng nghe nhạc trực tuyến"}}}\\
\\
\hline
\end{tabular}
\end{center}

\vspace{2cm}

\begin{table}[h]
\centering
    \begin{tabular}{rl}
    \hspace{3 cm}\textbf{GVHD}:
    & ThS. Trần Trương Tuấn Phát\\

    & \\[10pt]
\textbf{Sinh viên}: & Lư Chấn Vũ - 2313955 \emph{(Nhóm 10, \textbf{Leader})} \\
& Nguyễn Phú Vinh - 2313922 \emph{(Nhóm 10)} \\
& Trần Dương Khiết Nhi - 2312509 \emph{(Nhóm 10)} \\
& Lê Minh Khoa - 2311593 \emph{(Nhóm 10)} \\
& Lê Minh Trí - 2313593 \emph{(Nhóm 10)} \\
    \end{tabular}
\end{table}

\begin{center}
{\footnotesize TP. HỒ CHÍ MINH, 09/2025}
\end{center}
\end{titlepage}

\pagebreak
\tableofcontents

\pagebreak

% Glossaries
% {}
\printunsrtglossary[type={symbols}, title={Danh sách kí hiệu}]
\printunsrtglossary[type={abbreviations}, title={Danh sách từ viết tắt}]
\pagebreak
\listoffigures
\listoftables
\pagebreak
\addcontentsline{toc}{section}{\listfigurename}
\addcontentsline{toc}{section}{\listtablename}

% 

% Member list
\section*{Danh sách thành viên và nhiệm vụ}
\addcontentsline{toc}{section}{Danh sách thành viên và nhiệm vụ}
\begin{center}
\begin{table}[H]
\centering
\begin{tabular}{|c|c|c|l|c|}
\hline
\textbf{STT} & \textbf{Họ và tên} & \textbf{MSSV} & \textbf{Nhiệm vụ} & \textbf{\% hoàn thành}\\
\hline 
%%%%%Student 1%%%%%%%%%%
\multirow{3}{*}{1} & \multirow{3}{*}{Lư Chấn Vũ} & \multirow{3}{*}{2313955} & 
- & \multirow{3}{*}{100\%}\\
 & & & - & \\
\hline
%%%%%Student 2%%%%%%%%%%
\multirow{3}{*}{2} & \multirow{3}{*}{Nguyễn Phú Vinh} & \multirow{3}{*}{2313922} & 
- & \multirow{3}{*}{100\%}\\
 & & & - & \\
\hline
%%%%%Student 3%%%%%%%%%%
\multirow{3}{*}{3} & \multirow{3}{*}{Trần Dương Khiết Nhi} & \multirow{3}{*}{2312509} & 
-& \multirow{3}{*}{100\%}\\
 & & & - & \\
\hline
%%%%%Student 4%%%%%%%%%%
\multirow{3}{*}{4} & \multirow{3}{*}{Lê Minh Khoa} & \multirow{3}{*}{2311593} & 
- & \multirow{3}{*}{100\%}\\
 & & & - & \\
\hline
%%%%%Student 5%%%%%%%%%%
\multirow{3}{*}{5} & \multirow{3}{*}{Lê Minh Trí} & \multirow{3}{*}{2313593} & 
- & \multirow{3}{*}{100\%}\\
 & & & - & \\
\hline
\end{tabular}
\caption{\label{table1}Danh sách thành viên và nhiệm vụ}
\end{table}
\end{center}

\pagebreak
\section*{Nội dung báo cáo}

% Đăng nhập & Đăng ký - Người dùng
\begin{table}[h!]
\centering
\renewcommand{\arraystretch}{1.3} % tăng khoảng cách dòng trong bảng
\begin{tabularx}{\textwidth}{|l|X|}
\hline
\textbf{Use case name} & Đăng nhập \& Đăng ký \\ \hline
\textbf{Actors}        & Người dùng, Nghệ sĩ \\ \hline
\textbf{Description}   & Người Dùng có thể tạo tài khoản mới (Đăng ký) hoặc truy cập tài khoản đã có (Đăng nhập) để sử dụng ứng dụng. \\ \hline
\textbf{Trigger}       & Người dùng mở web và chọn tính năng “Đăng nhập / Đăng ký”. \\ \hline
\textbf{Pre-Condition(s)} 
& 1. Thiết bị của Người dùng phải được kết nối internet. \newline
  2. Người Dùng chưa có tài khoản (nếu đăng ký) hoặc đã có tài khoản hợp lệ (nếu đăng nhập). \\ \hline
\textbf{Post-Condition(s)} 
& Người Dùng đăng nhập thành công hoặc tạo mới tài khoản thành công và được chuyển vào giao diện chính, tài khoản đăng ký thành công sẽ được thêm vào hệ thống. \\ \hline
\textbf{Normal Flow}   
& \textbf{Đăng nhập:} \newline
  1. Người Dùng chọn “Đăng nhập”. \newline
  2. Nhập email/số điện thoại/tên đăng nhập và mật khẩu. \newline
  3. Hệ thống kiểm tra thông tin xác thực. \newline
  4. Nếu đúng $\rightarrow$ Đăng nhập thành công và hiển thị trang chủ. \newline
  \textbf{Đăng ký:} \newline
  1. Người Dùng chọn “Đăng ký”. \newline
  2. Nhập thông tin cá nhân: email/số điện thoại, mật khẩu, tên hiển thị. \newline
  3. Xác thực tài khoản qua email. \newline
  4. Hệ thống tạo tài khoản mới và hiển thị thông báo thành công. \\ \hline
\textbf{Exception Flow} 
& \textbf{Đăng nhập:} \newline
  2a. Bỏ trống trường thông tin $\rightarrow$ báo lỗi “Vui lòng nhập đầy đủ thông tin”. \newline
  3a. Sai mật khẩu $\rightarrow$ báo lỗi “Sai mật khẩu”. \newline
  3b. Tài khoản bị khóa $\rightarrow$ báo lỗi “Tài khoản của bạn đã bị khóa”.
  \textbf{Đăng ký:} \newline
  2a. Email/số điện thoại đã tồn tại $\rightarrow$ báo lỗi “Tài khoản đã tồn tại”. \newline
  2b. Mật khẩu yếu/không hợp lệ $\rightarrow$ báo lỗi “Mật khẩu không hợp lệ”. \newline
  3a. Không xác thực được email $\rightarrow$ báo lỗi “Xác thực thất bại”. \\ \hline
\textbf{Alternative Flow} 
& \textbf{Đăng nhập:} \newline
  2b. Người Dùng chọn “Đăng nhập bằng Google/Facebook/Apple ID”. \newline
  2c. Người Dùng chọn “Quên mật khẩu” $\rightarrow$ chuyển sang luồng khôi phục mật khẩu. \newline
  \textbf{Đăng ký:} \newline
  2c. Người Dùng có thể tải ảnh đại diện hoặc chọn sở thích âm nhạc ngay khi tạo tài khoản. \\ \hline
\end{tabularx}
\caption{Mô tả usecase Đăng nhập \& Đăng ký}
\end{table}

% Quản lý tài khoản - Người dùng
\begin{table}[h!]
\centering
\renewcommand{\arraystretch}{1.3} % tăng khoảng cách dòng trong bảng
\begin{tabularx}{\textwidth}{|l|X|}
\hline
\textbf{Use case name} & Quản lý tài khoản \\ \hline
\textbf{Actors}        & Người dùng, Nghệ sĩ \\ \hline
\textbf{Description}   & Người Dùng có thể quản lý tài khoản cá nhân bao gồm: đăng ký, đổi mật khẩu, cập nhật thông tin cá nhân (avatar, tên hiển thị, email), và xóa tài khoản. \\ \hline
\textbf{Trigger}       & Người dùng chọn chức năng “Quản lý tài khoản” trên giao diện hệ thống sau khi đăng nhập thành công. \\ \hline
\textbf{Pre-Condition(s)} 
& 1. Thiết bị của Người dùng phải được kết nối internet. \newline
  2. Người dùng đã đăng nhập thành công. \\ \hline
\textbf{Post-Condition(s)} 
& Hệ thống cập nhật thông tin tài khoản theo thao tác của Người Dùng. \\ \hline
\textbf{Normal Flow}   
& 1. Người dùng chọn vào phần “Quản lý tài khoản” trong giao diện hệ thống. \newline
  2. Hệ thống hiển thị thông tin tài khoản hiện tại. \newline
  3. Người dùng chọn hành động: chỉnh sửa thông tin cá nhân / đổi mật khẩu / thiết lập bảo mật. \newline
  4. Người dùng nhập thông tin mới hoặc thay đổi cần thiết. \newline
  5. Người dùng xác nhận và lưu thay đổi. \newline
  6. Hệ thống cập nhật dữ liệu và thông báo thành công. \\ \hline
\textbf{Exception Flow} 
& 3a. Nếu hệ thống không tải được thông tin tài khoản thì hiển thị lỗi “Không thể tải dữ liệu”. \newline
  4a. Nếu thông tin nhập sai định dạng (ví dụ email không hợp lệ, mật khẩu quá ngắn), hệ thống thông báo lỗi và yêu cầu nhập lại. \newline
  5a. Nếu kết nối internet bị gián đoạn trong khi lưu, hệ thống hiển thị thông báo thất bại và yêu cầu thử lại. \\ \hline
\textbf{Alternative Flow} 
& 3b. Người dùng có thể hủy thao tác và quay lại trang chính mà không thay đổi gì. \newline
  4b. Người dùng có thể bật/tắt các tính năng nâng cao như xác thực 2 lớp. \\ \hline
\end{tabularx}
\caption{Mô tả usecase Quản lý tài khoản}
\end{table}

% Tương tác - Người dùng
\begin{table}[h!]
\centering
\renewcommand{\arraystretch}{1.3} % tăng khoảng cách dòng trong bảng
\begin{tabularx}{\textwidth}{|l|X|}
\hline
\textbf{Use case name} & Tương tác \\ \hline
\textbf{Actors}        & Người dùng \\ \hline
\textbf{Description}   & Người dùng có thể thực hiện các hành động tương tác trong hệ thống như: thích (like) bài hát, bình luận bài hát, theo dõi nghệ sĩ, chia sẻ bài hát, nghệ sĩ, playlist lên mạng xã hội. \\ \hline
\textbf{Trigger}       & Người dùng chọn một bài hát, nghệ sĩ hoặc playlist bất kỳ trong hệ thống và mở giao diện chi tiết. \\ \hline
\textbf{Pre-Condition(s)} 
& 1. Thiết bị của Người dùng phải được kết nối internet. \newline
  2. Người dùng đã đăng nhập thành công. \newline
  3. Bài hát, nghệ sĩ đã có sẵn trên hệ thống. \\ \hline
\textbf{Post-Condition(s)} 
& Hệ thống lưu lại thông tin tương tác (like, bình luận, theo dõi, chia sẻ) và cập nhật dữ liệu liên quan. \\ \hline
\textbf{Normal Flow}   
& 1. Người dùng chọn một bài hát, nghệ sĩ, playlist. \newline
  2. Hệ thống hiển thị chi tiết bài hát, nghệ sĩ, playlist. \newline
  3. Người dùng chọn hành động muốn thực hiện: \newline
  \cach 3a. Nhấn nút “Thích” (Like) để thêm vào danh sách yêu thích. \newline
  \cach 3b. Viết và gửi bình luận cho bài hát. \newline
  \cach 3c. Nhấn “Theo dõi” để theo dõi nghệ sĩ. \newline
  \cach 3d. Chọn “Chia sẻ” và lựa chọn kênh chia sẻ (Các trang mạng xã \cach hội). \newline
  4. Hệ thống xác nhận và cập nhật thông tin tương tác. \\ \hline
\textbf{Exception Flow} 
& 2a. Nếu bài hát, nghệ sĩ, playlist không tồn tại hoặc bị xóa $\rightarrow$ hiển thị thông báo lỗi. \newline
  3b1. Nếu bình luận chứa nội dung vi phạm chính sách $\rightarrow$ hệ thống từ chối đăng và hiển thị thông báo. \newline
  3d1. Nếu chia sẻ thất bại do mất kết nối internet $\rightarrow$ hiển thị thông báo “Chia sẻ không thành công, vui lòng thử lại”. \\ \hline
\textbf{Alternative Flow} 
& 3a1. Người dùng có thể “Bỏ thích” (Unlike) nếu trước đó đã thích. \newline
  3b2. Người dùng có thể chỉnh sửa hoặc xóa bình luận đã đăng. \newline
  3c1. Người dùng có thể hủy theo dõi nghệ sĩ. \newline
  3d2. Người dùng có thể sao chép đường dẫn (link) thay vì chia sẻ trực tiếp. \\ \hline
\end{tabularx}
\caption{Mô tả usecase Tương tác}
\end{table}

% Nghe nhạc - Người dùng
\begin{table}[h!]
\centering
\renewcommand{\arraystretch}{1.3} % tăng khoảng cách dòng trong bảng
\begin{tabularx}{\textwidth}{|l|X|}
\hline
\textbf{Use case name} & Nghe nhạc \\ \hline
\textbf{Actors}        & Người dùng \\ \hline
\textbf{Description}   & Người dùng có thể phát và điều khiển nhạc với nhiều tính năng nâng cao: phát/tạm dừng, tiếp tục, chuyển tiếp/lùi lại, bật chế độ nghe ngẫu nhiên, nghe lặp lại (một bài hoặc toàn bộ danh sách), thêm bài vào hàng chờ, cài giờ tắt nhạc (sleep timer), và xem lời bài hát. \\ \hline
\textbf{Trigger}       & Người dùng chọn một bài hát, playlist, album và nhấn nút “Phát nhạc”. \\ \hline
\textbf{Pre-Condition(s)} 
& 1. Thiết bị của Người dùng phải được kết nối internet. \newline
  2. Ứng dụng được cấp quyền truy cập âm thanh. \newline
  3. Bài hát đã có sẵn trên hệ thống và có thể phát. \\ \hline
\textbf{Post-Condition(s)} 
& Hệ thống phát nhạc theo thao tác của Người dùng, cập nhật trạng thái trình phát (player state) và lưu lại lịch sử nghe nhạc. \\ \hline
\textbf{Normal Flow}   
& 1. Người dùng chọn một bài hát, playlist, album. \newline
  2. Hệ thống tải dữ liệu nhạc và bắt đầu phát. \newline
  3. Người dùng có thể thực hiện các thao tác điều khiển: \newline
  \cach 3a. Nhấn “Tạm dừng” (Pause) hoặc “Tiếp tục” (Play). \newline
  \cach 3b. Nhấn “Chuyển tiếp” (Next) hoặc “Lùi lại” (Previous). \newline
  \cach 3c. Bật chế độ “Nghe ngẫu nhiên” (Shuffle). \newline
  \cach 3d. Chọn chế độ “Nghe lặp lại” (Repeat one / Repeat all). \newline
  \cach 3e. Thêm bài hát vào hàng chờ phát (Queue). \newline
  \cach 3f. Cài đặt “Giờ tắt nhạc” (Sleep timer). \newline
  \cach 3g. Xem lời bài hát (Lyrics) nếu có sẵn. \newline
  4. Hệ thống phản hồi ngay lập tức và cập nhật trình phát nhạc. \newline
  5. Khi bài hát kết thúc, hệ thống tự động phát bài tiếp theo trong playlist/queue hoặc dừng lại nếu không còn bài nào \\ \hline
\textbf{Exception Flow} 
& 2a. Nếu bài hát không thể phát (do bản quyền hoặc lỗi file) $\rightarrow$  hiển thị thông báo. \newline
  2b. Nếu người dùng bị mất kết nối trong lúc đang phát nhạc $\rightarrow$ hiển thị thông báo “Mất kết nối, vui lòng kiểm tra lại” và tạm dừng phát nhạc. \newline
  3g1. Nếu bài hát không có lời (lyrics) trong cơ sở dữ liệu $\rightarrow$ hiển thị “Chưa có lời bài hát”. \\ \hline
\textbf{Alternative Flow} 
& 3a1. Người dùng có thể sử dụng tai nghe hoặc thiết bị ngoài để điều khiển (nút Play/Pause). \newline
  3e1. Người dùng có thể sắp xếp lại thứ tự bài hát trong hàng chờ. \newline
  3f2. Người dùng có thể hủy hoặc thay đổi thời gian sleep timer. \\ \hline
\end{tabularx}
\caption{Mô tả usecase Nghe nhạc}
\end{table}

% Khám phá và tìm kiếm - Người dùng
\begin{table}[h!]
\centering
\renewcommand{\arraystretch}{1.3} % tăng khoảng cách dòng trong bảng
\begin{tabularx}{\textwidth}{|l|X|}
\hline
\textbf{Use case name} & Khám phá và tìm kiếm \\ \hline
\textbf{Actors}        & Người dùng \\ \hline
\textbf{Description}   & Người dùng có thể tìm kiếm bài hát, nghệ sĩ, album, playlist và khám phá nhạc mới thông qua bảng xếp hạng, xu hướng (trending) và gợi ý cá nhân hóa. \\ \hline
\textbf{Trigger}       & Người dùng mở thanh tìm kiếm hoặc tab “Khám phá” trong ứng dụng. \\ \hline
\textbf{Pre-Condition(s)} 
& 1. Thiết bị của Người dùng phải được kết nối internet. \newline
  2. Dữ liệu về bài hát, nghệ sĩ, album phải có trong cơ sở dữ liệu. \\ \hline
\textbf{Post-Condition(s)} 
& Hệ thống trả về kết quả tìm kiếm hoặc danh sách nhạc khám phá phù hợp, cho phép Người dùng chọn và nghe nhạc ngay. \\ \hline
\textbf{Normal Flow}   
& 1. Người dùng mở tính năng Tìm kiếm/Khám phá. \newline
  2. Người dùng có thể thực hiện: \newline
  \cach 2a. Nhập từ khóa để tìm bài hát, nghệ sĩ, album, playlist. \newline
  \cach 2b. Xem danh sách Trending hoặc Top Chart. \newline
  \cach 2c. Nhận gợi ý nhạc cá nhân hóa dựa trên lịch sử nghe, lượt thích và nghệ sĩ theo dõi. \newline
  3. Hệ thống hiển thị danh sách kết quả. \newline
  4. Người dùng chọn nội dung mong muốn (phát nhạc, xem chi tiết nghệ sĩ, album, playlist). \\ \hline
\textbf{Exception Flow} 
& 2a1. Nếu không tìm thấy kết quả $\rightarrow$ hiển thị thông báo “Không tìm thấy nội dung phù hợp”. \newline
  2a2. Nếu không nhập gì $\rightarrow$ hiển thị “Vui lòng nhập từ khóa”. \newline
  2a3. Nếu kết nối internet bị gián đoạn trong khi tìm kiếm $\rightarrow$ hiển thị thông báo “Mất kết nối, vui lòng kiểm tra lại”. \newline
  2c1. Nếu Người dùng mới/ẩn danh (không đăng nhập) chưa có lịch sử nghe $\rightarrow$ gợi ý nhạc phổ biến mặc định. \\ \hline
\textbf{Alternative Flow} 
& 2a2. Người dùng có thể dùng bộ lọc nâng cao (theo thể loại, năm phát hành, nghệ sĩ, album…). \newline
  2b2. Người dùng có thể chọn xem bảng xếp hạng theo từng khu vực/quốc gia. \newline
  2c2. Người dùng có thể cập nhật gợi ý bằng cách thay đổi sở thích cá nhân (genres, mood…). \\ \hline
\end{tabularx}
\caption{Mô tả usecase Khám phá và tìm kiếm}
\end{table}

% Nhận diện bài hát - Người dùng
\begin{table}[h!]
\centering
\renewcommand{\arraystretch}{1.3} % tăng khoảng cách dòng trong bảng
\begin{tabularx}{\textwidth}{|l|X|}
\hline
\textbf{Use case name} & Nhận diện bài hát \\ \hline
\textbf{Actors}        & Người dùng \\ \hline
\textbf{Description}   & Người dùng bật tính năng nhận diện nhạc để hệ thống nghe một đoạn nhạc từ môi trường xung quanh, sau đó so sánh với cơ sở dữ liệu fingerprint để xác định bài hát. \\ \hline
\textbf{Trigger}       & Người dùng nhấn vào biểu tượng "Nhận diện bài hát" trên thanh tìm kiếm. \\ \hline
\textbf{Pre-Condition(s)} 
& 1. Thiết bị của Người dùng phải được kết nối internet. \newline
  2. Ứng dụng được cấp quyền truy cập micro. \newline
  3. Hệ thống có cơ sở dữ liệu nhạc đủ lớn để so khớp. \\ \hline
\textbf{Post-Condition(s)} 
& Hệ thống hiển thị thông tin bài hát (tên, nghệ sĩ, album) và cung cấp tùy chọn phát nhạc hoặc lưu vào playlist. \\ \hline
\textbf{Normal Flow}   
& 1. Người dùng chọn tính năng Nhận diện bài hát. \newline
  2. Hệ thống yêu cầu quyền truy cập micro và người dùng cấp phép. \newline
  3. Hệ thống ghi lại một đoạn âm thanh từ 15-20s. \newline
  4. Hệ thống trích xuất "Fingerprint" (dấu vân tay âm thanh) từ đoạn ghi âm để so khớp với dữ liệu. \newline
  5. Hệ thống trả về kết quả gồm tên bài hát, nghệ sĩ, album (nếu có). \newline
  6. Người dùng có thể chọn nghe bài hát, thêm vài playlist hoặc chia sẻ. \\ \hline
\textbf{Exception Flow} 
& 2a. Người dùng từ chối quyền micro $\rightarrow$ hệ thống báo không thể sử dụng tính năng. \newline
  3a. Không thu được âm thanh do quá ồn, micro hỏng $\rightarrow$ hệ thống yêu cầu thử lại. \newline
  5a. Nếu không tìm thấy kết quả khớp $\rightarrow$ hệ thống thông báo “Không nhận diện được bài hát”. \newline
  5b. Nếu có lỗi kết nối trong quá trình so khớp $\rightarrow$ hệ thống báo “Mất kết nối” và cho phép thử lại sau. \\ \hline
\textbf{Alternative Flow} 
& 5c. Người dùng có thể tìm kiếm bài hát thủ công. \\ \hline
\end{tabularx}
\caption{Mô tả usecase Nhận diện bài hát}
\end{table}
\end{document}